% 结构级反例样本：故意触发句长/句式重复/过渡词/段落节奏/小数位指标，仅用于检查器冒烟与对照；请勿照抄
% 对应指标：long-passive / repeated-opening / transition-density / sentence-rhythm / paragraph-rhythm / decimal-precision

\section{模型求解}

本文以最小化预测误差为目标建立梯度提升模型，并采用五折交叉验证选择超参数。本文构造滑动窗口滞后特征。本文以 TPE 搜索树深与学习率。本文最终得到的平均绝对百分比误差为 0.12345。

首先完成数据清洗，其次构造特征，然而特征量纲差异明显，因此需要标准化，最后输入模型训练。在上述处理下，本模型对早高峰流量具有较好的刻画能力，并且在多种扰动下均能保持稳定表现。受限于样本规模，模型参数被反复调整而未能收敛到全局最优，该现象在数据稀疏时段尤为明显。

本文认为该结果主要归因于特征工程的改进，并且这一判断得到交叉验证的支持。本文进一步分析了不同时段的误差分布，发现高峰时段的误差明显高于平峰时段。本文据此提出了分时段建模的改进思路。本文同时补充了特征重要性分析以解释模型行为。

模型在训练集与测试集上的表现较为接近，说明不存在严重过拟合。残差近似服从零均值正态分布，且未观察到明显的异方差结构。为验证稳健性，本文对学习率与树深分别作了扰动，结果显示指标波动较小。综合来看，该模型可用于信号配时的滚动评估。

结论表明该模型在给定样本范围内有效，但推广到其他城市仍需重新标定参数。本文给出了适用边界与后续改进方向。
